\documentclass[a4paper,11pt]{article}
\usepackage{sbpo-template}
\usepackage[brazil]{babel}    % português do Brasil
\usepackage[utf8]{inputenc}   % acentos em UTF-8
\usepackage[T1]{fontenc}      % para hifenização correta
\usepackage{amsmath,amssymb}
\usepackage{url}
\usepackage[square]{natbib}
\usepackage{indentfirst}
\usepackage{fancyhdr}
\usepackage{graphicx}
\usepackage{booktabs}
\usepackage{makecell}
\pagestyle{fancy}
\fancyhf{}
\fancyhead[C]{
\includegraphics[scale=0.5]{figuras/sbpo2026-header-logo.png}}
%\renewcommand{\headrulewidth}{0pt}
\renewcommand{\headruleskip}{-1mm}
%\setlength\headheight{101.0pt}
%\addtolength{\textheight}{-101.0pt}
\setlength\headheight{86pt}
\addtolength{\textheight}{-86pt}
\setlength{\headsep}{5mm}
\setlength{\footskip}{4.08003pt}

\title{Otimização de Inventário Multi-Echelon sob Demanda Intermitente: Comparação de Doze Políticas em Dados Reais de Varejo Brasileiro}

\author{
\name{}
\institute{Filiação}
\iaddress{Endereço da Instituição}
\email{e-mail}
}

\author{
\name{Segundo Autor}
\institute{Filiação}
\iaddress{Endereço da Instituição}
\email{e-mail}
}

\begin{document}

\maketitle

\vspace{8mm}
\begin{resumo}
Modelos clássicos de inventário, como o \textit{Economic Order Quantity}
(EOQ) e a política de dois limiares~$(s,S)$, pressupõem
estacionariedade, hipótese que apresenta limitações sistemáticas
em cenários com demanda intermitente, gerando excesso de estoque
e rupturas frequentes no varejo regional brasileiro. Estudos comparativos recentes baseados em
conjuntos como o Walmart~M5 não cobrem o regime de alta variabilidade
caracterizado por coeficiente de variação (CV) entre $2.5$ e $5.4$,
típico de produtos de baixo giro. Este trabalho apresenta uma
comparação de doze políticas de controle de inventário em dados
operacionais proprietários de varejo regional brasileiro, avaliadas
em cinco indicadores: custo total, nível de serviço, taxa de ruptura,
efeito de amplificação da demanda (\textit{bullwhip}) e frequência
de pedidos. A arquitetura híbrida GA-DQN
(\textit{Genetic Algorithm with Deep Q-Network}) apresentou o
equilíbrio mais favorável entre as políticas avaliadas: custo
inferior às heurísticas clássicas, efeito \textit{bullwhip} menor
que a política~$(s,S)$ e taxa de ruptura equivalente à melhor
política prática isolada, resultado consistente com o desempenho
esperado da arquitetura em regime de alta intermitência.
\end{resumo}


\bigskip
\begin{palchaves}
Gestão de inventário; Demanda intermitente; Cadeia de suprimentos
multi-nível; Meta-heurísticas; Aprendizado por reforço.

\bigskip
\noindent{Tópicos (em ordem de PRIORIDADE): }
Logística e \textit{Supply Chain}; Otimização de Processos;
Inteligência Artificial e Aprendizado de Máquina.
\end{palchaves}

\vspace{8mm}
\begin{abstract}
Classical inventory models, such as the Economic Order Quantity (EOQ)
and the two-threshold policy~$(s,S)$, rely on stationarity assumptions, a hypothesis that presents
systematic limitations for intermittent-demand items,
producing excess stock and frequent stockouts in Brazilian regional
retail.
Recent comparative studies based on datasets such as Walmart~M5 do
not cover this high-variability regime, characterized by a coefficient
of variation (CV) between $2.5$ and $5.4$, typical of slow-moving
products. This work presents a comparison of twelve inventory control policies
on proprietary operational data from Brazilian regional retail,
evaluated on five indicators: total cost, service level, stockout
rate, demand amplification along the supply chain (bullwhip effect),
and order frequency. The hybrid GA-DQN (Genetic Algorithm with Deep
Q-Network) architecture presented the most favorable trade-off among
the evaluated policies: lower cost than classical heuristics, a
bullwhip effect lower than $(s,S)$, and a stockout rate on par with
the best practical standalone policy, suggesting consistent
performance of the hybrid approach under high intermittent demand.
\end{abstract}

\bigskip
\begin{keywords}
Inventory management; Intermittent demand; Multi-level supply chain;
Metaheuristics; Reinforcement learning.

\bigskip
\noindent{Paper topics (in order of PRIORITY): }
Logistics and Supply Chain; Process Optimisation;
Artificial Intelligence and Machine Learning.
\end{keywords}

% =========================================================
\section{Introdução}
% =========================================================

A gestão ineficiente de inventário representa, segundo estimativas,
entre $3.5\%$ e $8\%$ do faturamento anual do varejo brasileiro em
capital imobilizado e perdas por ruptura~\citep{silver:2017}. No
varejo regional, itens de baixo giro são tipicamente controlados com
heurísticas clássicas como o Lote Econômico de Compra
(EOQ)~\citep{harris:1913} e a política~$(s,S)$~\citep{hadley:1963},
que pressupõem estacionariedade raramente verificada em dados
reais~\citep{silver:2017}. Em cadeias \textit{multi-echelon} com
demanda \emph{intermitente}, longos períodos sem consumo alternados
com picos esporádicos~\citep{croston:1972} amplificam esse problema:
violam hipóteses de normalidade e intensificam o efeito
\textit{bullwhip} ao longo da cadeia. Estudos comparativos recentes
baseados no Walmart~M5 não cobrem o regime de coeficiente de
variação (CV) superior a $2.5$, típico de produtos de baixo giro
no varejo regional brasileiro~\citep{syntetos:2001}.

Apesar dos avanços recentes em RL e meta-heurísticas, há pouca
evidência empírica comparando essas abordagens em cadeias
\textit{multi-echelon} sob demanda altamente intermitente baseada
em dados operacionais proprietários de varejo regional brasileiro,
especialmente em regime de CV entre $2.5$ e $5.4$.
A hipótese central é que arquiteturas híbridas
GA-RL (\textit{Genetic Algorithm -- Reinforcement Learning})
apresentam equilíbrio mais favorável entre custo, estabilidade
logística e nível de serviço do que políticas clássicas e agentes
de RL isolados nesse contexto. As contribuições deste trabalho são:

\begin{enumerate}
  \item Comparação sistemática de doze políticas de controle
    de inventário em dados operacionais proprietários de varejo
    regional brasileiro com $CV$ entre $2.5$ e $5.4$, regime
    pouco explorado na literatura disponível;
  \item Avaliação em regime de demanda extremamente intermitente,
    não coberto por estudos baseados no Walmart~M5;
  \item Análise com cinco indicadores primários para todas as
    políticas, incluindo efeito \textit{bullwhip} e frequência
    de pedidos;
  \item Estratificação por perfil de demanda, revelando que a
    política ótima depende do regime de consumo;
  \item Validação da arquitetura GA-DQN
    (\textit{Genetic Algorithm -- Deep Q-Network}) em dados
    operacionais proprietários de varejo regional brasileiro.
\end{enumerate}

Os resultados sugerem que o GA-DQN apresentou o equilíbrio mais
favorável entre as políticas avaliadas: custo $48\%$ inferior à
política~$(s,S)$, efeito \textit{bullwhip} aproximadamente 24
vezes menor e taxa de ruptura equivalente à melhor política
isolada. A política~$(s,S)$, frequentemente utilizada no varejo
regional, apresentou desempenho consistentemente inferior em todos
os estratos avaliados.

O restante do artigo está organizado como se segue: a Seção~2
revisa a literatura; a Seção~3 apresenta a formulação do problema;
a Seção~4 descreve os dados; a Seção~5 detalha as doze políticas;
a Seção~6 descreve o ambiente experimental; a Seção~7 apresenta e
discute os resultados; e a Seção~8 conclui com recomendações
práticas e trabalhos futuros.

% =========================================================
\section{Trabalhos Relacionados}
% =========================================================

Esta seção revisa os trabalhos mais relacionados em três frentes:
o problema de inventário \textit{multi-echelon} sob demanda
intermitente, as abordagens de meta-heurísticas e aprendizado por
reforço aplicadas ao controle de estoque, e as lacunas da literatura
que este trabalho endereça.

\subsection{Inventário Multi-Echelon sob Demanda Intermitente}

A teoria clássica de controle de estoques, fundamentada no
EOQ~\citep{harris:1913} e na política~$(s,S)$~\citep{hadley:1963},
pressupõe demanda estacionária e quase-normal. Em contextos de
demanda \emph{intermitente}, essa hipótese apresenta limitações
sistemáticas: a distribuição empírica é bimodal (zero com alta
probabilidade; pico raro de grande magnitude) e os estimadores
clássicos perdem confiabilidade~\citep{croston:1972}. \citet{syntetos:2001}
formalizaram o regime intermitente em termos do coeficiente de
variação, estabelecendo que padrões com CV~$>1.5$ exigem tratamento
metodológico específico, limiar superado por todos os cenários
deste estudo.

No ambiente \textit{multi-echelon}, a complexidade adicional decorre
do acoplamento entre níveis: a variabilidade dos pedidos das lojas
amplifica-se ao longo da cadeia (efeito
\textit{bullwhip}~\citep{lee:1997}), tornando o problema de
reposição ainda mais difícil de resolver com parâmetros
estáticos~\citep{axsater:2006,clark:1960}. Problemas de grande
escala, como o conjunto completo que origina este estudo (48 milhões
de transações e 12 mil SKUs), evidenciam que o espaço
de decisões cresce combinatorialmente com o número de produtos,
lojas e ciclos~\citep{shapiro:2001,boute:2022}.

\subsection{Meta-Heurísticas e Aprendizado por Reforço}

Uma primeira frente de pesquisa emprega meta-heurísticas para
otimização de parâmetros de inventário. \citet{yuna:2023}
compararam \textit{Simulated Annealing} (SA) e Busca Tabu na
otimização de parâmetros~$(s,S)$ em sete conjuntos de dados
intermitentes, verificando superioridade do SA em cinco dos sete
cenários. \citet{erkayman:2025} modelaram a chegada de demanda como
cadeia de Markov e otimizaram parâmetros de reposição com GA e PSO
(\textit{Particle Swarm Optimization}), mostrando que o PSO
converge mais rapidamente em espaços de baixa dimensão. Ambos os
estudos, porém, restringem-se a um único escalão e não incorporam
agentes de RL.

Uma segunda frente aplica aprendizado por reforço ao controle de
inventário. \citet{oroojlooyjadid:2022} demonstraram, no \textit{Beer
Game} \textit{multi-echelon}, que DQN (\textit{Deep Q-Network}) aprende
políticas de reposição sem modelar a distribuição de demanda.
\citet{gijsbrechts:2022} e \citet{vanvuchelen:2020} confirmaram que o
RL profundo pode igualar, mas não domina universalmente, as melhores
heurísticas. \citet{liu:2024} estenderam o paradigma ao ambiente
descentralizado com HAPPO (\textit{Heterogeneous-Agent Proximal Policy
Optimization}), reduzindo custo e efeito \textit{bullwhip}
simultaneamente. Uma limitação transversal é a exigência de da ordem
de $10^6$ interações para convergência, tornando esses agentes
inadequados para séries históricas curtas.

\citet{zabraoui:2025} constituem a referência mais próxima deste
estudo. Sobre o conjunto Walmart~M5, os autores propuseram a
arquitetura híbrida GA-DQN: o GA otimiza globalmente os parâmetros
estáticos de inventário enquanto a DQN refina localmente as decisões
a cada ciclo. Essa combinação elevou o nível de serviço de $61\%$
(DQN isolado) para $94\%$, com redução simultânea do custo total.
O presente trabalho parte desse resultado e o estende a um regime
de CV entre $2.5$ e $5.4$, incorporando sete políticas adicionais
e cinco indicadores primários.

\subsection{Lacunas da Literatura}

A revisão acima revela quatro lacunas que este trabalho endereça:
(i)~ausência de estudos com CV extremo ($>2.5$) em dados reais de
varejo regional brasileiro; (ii)~falta de comparação integrada entre
políticas clássicas, meta-heurísticas e RL em ambiente
\textit{multi-echelon}; (iii)~avaliação restrita de indicadores:
efeito \textit{bullwhip} e taxa de ruptura raramente aparecem como
critérios primários para todas as políticas; e (iv)~ausência de
validação de arquiteturas híbridas fora do regime Walmart~M5.

A Tabela~\ref{tab:relwork} sintetiza as dimensões em que este
estudo avança o estado da arte.

\begin{table}[ht]
\centering
\small
\caption{Comparação de trabalhos relacionados em gestão de inventário.
\textbf{Negrito}: dimensões em que este trabalho avança o estado da arte.
``Dados Reais'': dados proprietários coletados no domínio-alvo
(varejo com alta intermitência); Walmart~M5, embora real, não
representa esse regime (CV típico $<1$).}
\label{tab:relwork}
\begin{tabular}{@{}lccccc@{}}
\toprule
Trabalho &
\makecell[c]{Dados\\Reais} &
\makecell[c]{CV\\(regime)} &
\makecell[c]{Políticas\\(\#)} &
\makecell[c]{Multi-\\nível} &
\makecell[c]{Indicadores\\(\#)} \\
\midrule
\citealt{oroojlooyjadid:2022} & Não & $<1$         &  3 & Sim & 2 \\
\citealt{vanvuchelen:2020}    & Não & $<1$         &  4 & Não & 1 \\
\citealt{gijsbrechts:2022}    & Não & $<1$         &  6 & Sim & 2 \\
\citealt{yuna:2023}           & Não & $>1.5$       &  2 & Não & 2 \\
\citealt{erkayman:2025}       & Não & $>1.5$       &  2 & Não & 1 \\
\citealt{liu:2024}            & Não & $<1$         &  4 & Sim & 2 \\
\citealt{zabraoui:2025}       & Não & $<1$         &  7 & Sim & 3 \\
\midrule
\textbf{Este trabalho}
  & \textbf{Sim}
  & $\mathbf{2.5}$--$\mathbf{5.4}$
  & \textbf{12}
  & \textbf{Sim}
  & \textbf{5} \\
\bottomrule
\end{tabular}
\end{table}

% =========================================================
\section{Formulação do Problema}
% =========================================================

\subsection{Definição Formal do Problema MEIO}

O problema \textit{Multi Echelon Inventory Optimization} (MEIO)
considera uma cadeia de suprimentos real com $N = 15$ lojas
varejistas abastecidas por um único armazém regional. O horizonte é
discretizado em ciclos bimestrais $t \in \{1, \ldots, T\}$, com
$T = 38$. As variáveis e parâmetros centrais são:

\begin{itemize}
  \item $D_{t,i} \geq 0$: demanda observada na loja $i$ no ciclo $t$;
  \item $I_{t,i}$: nível de estoque da loja $i$ no início de $t$;
  \item $Q_{t,i} \geq 0$: quantidade pedida em $t$, recebida após
    prazo $L = 2$ ciclos;
  \item $S_{t,i} = \max(0, D_{t,i} - I_{t,i})$: ruptura no ciclo $t$;
  \item $s_i$: ponto de reposição, acionado quando $I_{t,i} \leq s_i$;
  \item $h$, $c_s$, $c_o^{fix}$, $c_o^{var}$: custos unitários de
    manutenção, ruptura, ordenação fixa e variável.
\end{itemize}

Uma \emph{política} $\pi$ mapeia o estado observado
$(I_{t,i},\, D_{t-1,i},\, Q_{t-L,i})$ em uma decisão de pedido
$Q_{t,i} = \pi(\cdot)$. O objetivo é determinar $\pi^*$ que minimize
o custo total acumulado:
\begin{equation}
  \min_\pi \sum_{t=1}^{T} \bigl[
    h\, I_{t,i} + c_s\, S_{t,i}
    + c_o^{fix}\,\mathbf{1}(Q_{t,i}>0)
    + c_o^{var}\, Q_{t,i}
  \bigr],
  \label{eq:obj}
\end{equation}
sujeito a:
\begin{align}
  I_{t+1,i} &= \max\bigl(0,\; I_{t,i} - D_{t,i}\bigr) + Q_{t-L,i},
    & \forall\, t,\; \forall\, i, \label{eq:balance} \\
  \sum_{i=1}^{N} Q_{t,i} &\leq W_t,
    & \forall\, t, \label{eq:avail} \\
  I_{1,i} &= I_0,
    & \forall\, i, \label{eq:init} \\
  Q_{t,i} &\geq 0,
    & \forall\, t,\; \forall\, i. \label{eq:nonneg}
\end{align}
A restrição~\eqref{eq:balance} governa a dinâmica de estoque: pedidos
emitidos no ciclo $t$ chegam após prazo fixo $L = 2$ ciclos.
A restrição~\eqref{eq:avail} limita o total pedido pelas $N$ lojas
à disponibilidade $W_t$ do armazém no ciclo $t$; neste estudo,
assume-se disponibilidade irrestrita ($W_t \to \infty$),
simplificação representativa de redes em que o armazém não é o
gargalo operacional. A condição~\eqref{eq:init} fixa o estoque
inicial, escalonado por loja conforme a Seção~6.2. A
restrição~\eqref{eq:nonneg} impede pedidos negativos; o estoque
$I_{t,i}$ é naturalmente não-negativo pela dinâmica $\max(0,\cdot)$
em~\eqref{eq:balance}.

O regime de demanda é quantificado pelo coeficiente de variação
$CV = \hat{\sigma}_D / \hat{\mu}_D$, onde $\hat{\sigma}_D$ é o
desvio padrão e $\hat{\mu}_D$ é a média da demanda. Valores acima
de \(0.5\) indicam demanda errática e acima de \(1.5\), demanda
extremamente intermitente~\citep{syntetos:2001}. Todos os 15
cenários deste estudo possuem $CV$ entre \(2.5\) e \(5.4\), regime
em que as hipóteses clássicas apresentam limitações severas.

\subsection{Arquitetura Multi-Echelon}

A cadeia de suprimentos modela um sistema de inventário multi-nível
com três escalões hierárquicos, estrutura referenciada na literatura
como MEIO~\citep{axsater:2006,clark:1960}:

\begin{itemize}
  \item Escalão 1 (Armazém Regional): único centro distribuidor da
    Paraíba; repõe junto ao fornecedor com prazo $L = 2$ ciclos.
  \item Escalão 2 (Lojas Revendedoras): 15 unidades varejistas com
    séries de demanda independentes; cada loja decide individualmente
    quando e quanto pedir ao armazém.
  \item Escalão 3 (Produto/SKU): produto 48130, único SKU estudado;
    foco em um item isola o comportamento intermitente de efeitos de
    portfólio e torna os 38 ciclos suficientes para calibração e
    avaliação.
\end{itemize}

As decisões de reposição são tomadas de forma descentralizada: cada
política $\pi_i$ observa o estado local $(I_{t,i}, D_{t-1,i},
Q_{t-L,i})$ e emite o pedido $Q_{t,i}$ sem coordenação explícita
com as demais lojas. Essa arquitetura é representativa da realidade
operacional da maioria das redes de varejo regional brasileiras, onde
sistemas integrados de planejamento multi-loja raramente estão
disponíveis.

\subsection{Métricas de Avaliação}

Cinco indicadores são computados ao fim dos 38 ciclos para cada
combinação (política, loja):

\begin{itemize}
  \item \emph{CTI} (Custo Total de Inventário):
    $\sum_{t}[h I_t + c_s S_t + c_o^{fix}\mathbf{1}(Q_t>0) + c_o^{var}Q_t]$;
  \item \emph{NS} (Nível de Serviço):
    $\sum_t \min(I_t, D_t) / \sum_t D_t$;
  \item \emph{TR} (Taxa de Ruptura): fração de ciclos com \textit{stockout};
  \item \emph{BE} (\textit{Bullwhip Effect}):
    $\mathrm{Var}(Q_1,\ldots,Q_T)/\mathrm{Var}(D_1,\ldots,D_T)$,
    razão entre a variância dos pedidos e a variância da demanda,
    mede a amplificação da variabilidade ao longo da cadeia;
  \item \emph{FP} (Frequência de Pedidos): fração de ciclos com
    ao menos um pedido emitido.
\end{itemize}

% =========================================================
\section{Base de Dados e Pré-Processamento}
% =========================================================

Esta seção descreve o contexto e a escala do conjunto de dados
original, os critérios de seleção de lojas e produto, o
\textit{pipeline} de pré-processamento e a caracterização
estatística da demanda das quinze unidades experimentais.

\subsection{Contexto dos Dados}

Os dados provêm de um sistema transacional de gestão de vendas de
uma rede de varejo do Nordeste brasileiro. O conjunto completo
contém aproximadamente 48 milhões de registros de transações de
venda, distribuídos entre cerca de 12 mil produtos únicos (SKUs) e
mais de 3 milhões de entidades de demanda em múltiplos estados.
Para viabilizar o estudo experimental com rigor metodológico,
adotou-se um recorte com três critérios: (i)~restrição geográfica
ao estado da Paraíba (PB); (ii)~restrição de produto a um único
SKU de comportamento intermitente representativo; e (iii)~restrição
temporal ao período disponível com séries suficientemente longas.

\subsection{Estrutura Temporal e Ciclos de Venda}

O sistema registra vendas por \emph{ciclo bimestral} (17 ciclos por
ano, $\approx 21$ dias corridos cada), granularidade que reflete a
política de consolidação de pedidos ao armazém e justifica o prazo
$L = 2$ ciclos ($\approx 42$ dias). O horizonte experimental
compreende janeiro de 2023 a abril de 2025 ($T = 38$ ciclos),
período suficiente para capturar sazonalidade anual e padrões de
intermitência.

\subsection{Seleção do SKU e das Lojas}

O produto 48130 foi selecionado a partir do conjunto completo de
12 mil SKUs com base em dois critérios: (i)~maior cobertura de
lojas com demanda ativa (ao menos 17 ciclos de demanda positiva no
horizonte estudado); e (ii)~regime de demanda extremamente
intermitente (CV~$>1.5$ em todas as lojas ativas). O produto 48130
atende simultaneamente aos dois critérios, com 33 lojas ativas na
Paraíba, sendo representativo de artigos sazonais de baixo giro
com padrão errático de consumo.

\subsection{Pipeline de Pré-Processamento}

O pré-processamento seguiu cinco etapas sequenciais:
(i)~extração e filtragem geográfica (estado~=~PB, produto~=~48130);
(ii)~limpeza de anomalias (remoção de transações com quantidade
negativa ou superior a 50.000 unidades);
(iii)~agregação por ciclo (soma de quantidades vendidas por loja e
ciclo);
(iv)~reconstituição de séries completas (vetor de comprimento
$T = 38$ com imputação de zero nos ciclos sem registro);
(v)~seleção das quinze lojas com ao menos 17 ciclos com demanda
positiva ($45\%$ do horizonte), limiar necessário para
calibração estatística das políticas.

\subsection{Caracterização Estatística da Demanda por Loja}

A Tabela~\ref{tab:demanda} apresenta o perfil das quinze lojas.
A heterogeneidade é expressiva: a demanda média varia de $11.7$
(Loja~15) a $1455.2$ unidades por ciclo (Loja~1), razão de 124
vezes. Nenhuma loja apresenta $CV$ inferior a $2.5$, situando todos
os cenários bem acima do limiar de demanda extremamente intermitente
(CV~$>1.5$) de \citet{syntetos:2001} e invalidando as hipóteses de
normalidade subjacentes a EOQ e~$(s,S)$.

Esse gradiente motivou a estratificação em três grupos e a
calibração adaptativa descrita na Seção~6:

\begin{itemize}
  \item Grupo I (5 lojas, $\hat{\mu} > 500$ unid./ciclo):
    alto volume com picos esparsos ($CV$ entre $2.83$ e $5.39$);
    estoque inicial adequado garante NS\,=\,$1.00$ para quase todas
    as políticas, tornando o custo o principal diferenciador.
  \item Grupo II (6 lojas, $50 < \hat{\mu} \leq 500$):
    estrato de maior relevância prática, com variabilidade entre
    $2.83$ e $4.51$ e volume suficiente para que as políticas se
    diferenciem em todas as cinco métricas.
  \item Grupo III (4 lojas, $\hat{\mu} \leq 50$):
    baixíssimo volume; o CV próximo de 3 cria uma barreira estrutural
    em que qualquer pedido compatível com a média pode exceder a
    demanda real de vários ciclos futuros.
\end{itemize}

\begin{table}[ht]
\centering
\small
\caption{Perfil de demanda das quinze lojas selecionadas, produto
48130, Paraíba, 2023--2025. Lojas ordenadas por $\hat{\mu}^+$
decrescente. $\hat{\mu}^+$: média apenas nos ciclos com demanda
positiva; critérios de grupo baseados na média incondicional
$\hat{\mu}$ (sobre todos os 38 ciclos):
I ($\hat{\mu}>500$), II (50--500), III ($\leq 50$).}
\label{tab:demanda}
\begin{tabular}{@{}lrrrrc@{}}
\toprule
\makecell[l]{Loja} &
\makecell[r]{$\hat{\mu}^+$\\(unid./ciclo)} &
\makecell[r]{Desv. Pad.} &
\makecell[r]{CV} &
\makecell[r]{Ciclos\\Positivos} &
\makecell[c]{Grupo} \\
\midrule
1  & 1.455{,}2 & 5.037{,}5 & 5{,}14 & 18 & I \\
2  & 1.375{,}7 & 4.464{,}6 & 4{,}35 & 22 & I \\
3  &   703{,}6 & 2.561{,}1 & 5{,}39 & 18 & I \\
4  &   553{,}4 & 1.380{,}2 & 3{,}67 & 19 & I \\
5  &   518{,}2 & 1.150{,}1 & 2{,}83 & 25 & I \\
6  &   300{,}9 &   590{,}3 & 2{,}95 & 19 & II \\
7  &   246{,}0 &   546{,}4 & 3{,}04 & 22 & II \\
8  &   181{,}6 &   393{,}0 & 2{,}90 & 23 & II \\
9  &   169{,}7 &   489{,}5 & 3{,}79 & 23 & II \\
10 &    88{,}3 &   274{,}8 & 4{,}51 & 19 & II \\
11 &    85{,}4 &   253{,}3 & 4{,}31 & 19 & II \\
12 &    44{,}8 &    85{,}9 & 2{,}53 & 24 & III \\
13 &    43{,}1 &    88{,}4 & 2{,}83 & 22 & III \\
14 &    34{,}7 &    97{,}1 & 3{,}97 & 20 & III \\
15 &    11{,}7 &    25{,}5 & 3{,}24 & 19 & III \\
\bottomrule
\end{tabular}
\end{table}

% =========================================================
\section{Políticas de Controle de Inventário}
% =========================================================

Esta seção apresenta as doze políticas organizadas em quatro
categorias de crescente sofisticação computacional.

\subsection{Políticas Clássicas}

As três políticas clássicas fundamentam-se em hipóteses de
estacionariedade da demanda~\citep{silver:2017}: o EOQ minimiza
os custos de ordenação e manutenção sob demanda constante~\citep{harris:1913};
a~$(s,S)$ emite pedido $S-I_t$ ao cruzar o ponto de reposição
$s$~\citep{hadley:1963}; e o Jornaleiro determina o pedido ótimo
de período único como um quantil da distribuição de
demanda~\citep{cachon:2009}. Todas apresentam limitações em regime
intermitente~\citep{silver:2017}.

\emph{EOQ.} $Q^* = \sqrt{2 D c_o^{fix} / h}$, com ponto de
reposição $s = \hat{\mu} L$~\citep{harris:1913}.

\emph{$(s,S)$.} Pedido de magnitude $S - I_t$ sempre que
$I_t \leq s$, com $s = s_\text{EOQ}$ e
$S = s + Q_\text{EOQ}$~\citep{hadley:1963}.

\emph{Jornaleiro.} $Q^* = \hat{\mu} + z\hat{\sigma}$,
$z = 1.28$ (nível de serviço alvo de $90\%$)~\citep{cachon:2009}.

\subsection{Meta-Heurísticas de Otimização}

As quatro meta-heurísticas buscam $(s^*, Q^*)$ minimizando a
Equação~\eqref{eq:obj} por busca global no espaço de parâmetros:
GA~\citep{holland:1975} por seleção, cruzamento e mutação;
SA~\citep{kirkpatrick:1983} por aceitação probabilística de soluções
piores com temperatura decrescente ($\exp(-\Delta C / T_k)$);
PSO~\citep{kennedy:1995} por combinação de experiências individual
e coletiva de um enxame; e ED~\citep{storn:1997} por perturbação
diferencial na estratégia \textit{rand/1/bin}.

\subsection{Agentes de Aprendizado por Reforço}

Os agentes de RL aprendem a política diretamente da interação com
o ambiente, sem especificação explícita de $(s, Q)$. O vetor de
estado compartilhado contém seis componentes normalizados: nível de
estoque atual, demanda esperada, pedidos em trânsito, média móvel e
desvio padrão dos últimos seis ciclos, e progresso do horizonte.
A ação é $Q_{t,i}$ e a recompensa é o negativo do custo incorrido.

A DQN~\citep{mnih:2015} utiliza \textit{experience replay} e
\textit{target network} para estabilizar o \textit{Q-learning};
sua limitação central é a exigência de da ordem de $10^6$ interações
para convergência~\citep{oroojlooyjadid:2022}, inadequada para os
38 ciclos deste estudo. O PPO~\citep{schulman:2017} treina política
estocástica com \textit{clipping} da razão de probabilidades em
$[1-\epsilon_\text{clip}, 1+\epsilon_\text{clip}]$, garantindo
atualizações estáveis em ações contínuas, mas sensível ao ponto de
partida. O SARSA~\citep{sutton:2018} é \textit{on-policy} tabular,
conservador em ambientes ruidosos, porém limitado em expressividade
e sujeito a convergência degenerada em horizontes curtos.

\subsection{Arquiteturas Híbridas}

Arquiteturas híbridas combinam a busca global das meta-heurísticas
com a adaptação local dos agentes de RL. \citet{zabraoui:2025}
formalizaram esse princípio na arquitetura GA-DQN, que este trabalho
replica e estende com a variante GA-PPO em regime de CV entre $2.5$
e $5.4$, não coberto por \citet{zabraoui:2025}.

O \textit{pipeline} opera em três fases: (i)~o GA executa 50
gerações e identifica parâmetros de alta qualidade;
(ii)~as melhores soluções pré-populam um \textit{replay buffer}
com 1.000 transições representativas; (iii)~o agente de RL
(DQN ou PPO) treina por 200 episódios a partir desse ponto
qualificado, evitando exploração do zero em horizontes curtos.

\subsection{Hiperparâmetros}

A Tabela~\ref{tab:params} consolida os hiperparâmetros de todas as
abordagens, facilitando reprodutibilidade.

\begin{table}[ht]
\centering
\small
\caption{Hiperparâmetros das doze políticas de controle de
inventário.}
\label{tab:params}
\begin{tabular}{@{}llp{5.5cm}@{}}
\toprule
Política & Categoria & Hiperparâmetros Principais \\
\midrule
EOQ
  & Clássica
  & $Q^* = \sqrt{2Dc_o^{fix}/h}$;\; $s = \hat{\mu}L$ \\
$(s,S)$
  & Clássica
  & $s = s_\text{EOQ}$;\; $S = s + Q_\text{EOQ}$ \\
Jornaleiro
  & Clássica
  & $Q^* = \hat{\mu} + 1.28\hat{\sigma}$;\; SL alvo = $90\%$ \\
\midrule
GA
  & Metaheurística
  & Pop.\ = 100;\; Ger.\ = 50;\; $p_c = 0.7$;\; torneio $k=3$;\; mut.\ gaussiana \\
SA
  & Metaheurística
  & $T_0 = 1000$;\; $\alpha = 0.95$;\; aceitação: $\exp(-\Delta C/T_k)$ \\
PSO
  & Metaheurística
  & $N = 40$ part.;\; 80 iter.;\; $w = 0.7$;\; $c_1 = c_2 = 1.5$ \\
ED
  & Metaheurística
  & \textit{rand/1/bin};\; $F = 0.8$;\; $C_R = 0.9$ \\
\midrule
DQN
  & Reforço
  & 2 camadas $\times$ 64 ReLU;\; 500 ep.;\; $\epsilon$-greedy decrescente;\; \textit{replay} 10.000 \\
PPO
  & Reforço
  & 500 ep.;\; $\epsilon_\text{clip} = 0.2$;\; ações contínuas \\
SARSA
  & Reforço
  & Tabela-$Q$ 5$\times$10;\; $\alpha = 0.1$;\; $\gamma = 0.99$;\; $\epsilon = 0.1$ \\
\midrule
GA-DQN
  & Híbrida
  & GA (50 ger.) $\to$ \textit{buffer} 1.000 trans.\ $\to$ DQN 200 ep. \\
GA-PPO
  & Híbrida
  & GA (50 ger.) $\to$ \textit{buffer} 1.000 trans.\ $\to$ PPO 200 ep. \\
\bottomrule
\end{tabular}
\end{table}

% =========================================================
\section{Ambiente Experimental}
% =========================================================

Esta seção descreve o ambiente de simulação, a estratégia de
calibração adaptativa que garante comparação justa entre lojas de
volumes muito distintos e o protocolo de comparação adotado.

\subsection{Ambiente de Simulação}

O ambiente \textit{InventoryEnv}, implementado em Python~3.11,
simula um sistema de inventário $(s, Q)$ com prazo de entrega fixo
$L = 2$ ciclos. As bibliotecas utilizadas são NumPy, Pandas e
Stable-Baselines3 com \textit{backend} PyTorch para os agentes de
RL. Todas as simulações utilizam semente aleatória fixa
(\textit{seed}~=~42) para garantir reprodutibilidade. A série de
demanda real de cada loja é mantida idêntica para todas as políticas,
de modo que as diferenças de desempenho reflitam exclusivamente as
características de cada política.

\subsection{Calibração Adaptativa de Parâmetros por Loja}

A heterogeneidade entre as lojas ($11.7$ a $1455$ unid./ciclo)
inviabiliza o uso de parâmetros absolutos uniformes. Os parâmetros
de referência são escalonados pelas estatísticas $\hat{\mu}$ e
$\hat{\sigma}$ de cada loja:
\begin{align}
  I_0 &= \max\!\bigl(100,\; 2\hat{\mu} L\bigr), \label{eq:I0}\\
  Q_\text{ref} &= \sqrt{2\hat{\mu} T\, c_o^{fix} / h}, \label{eq:qref}\\
  s_\text{ref} &= \hat{\mu} L + z\hat{\sigma}\sqrt{L}, \label{eq:sref}
\end{align}
com $z = 1.28$ (nível de serviço alvo de $90\%$). Os espaços de
busca das meta-heurísticas e os limites de ação dos agentes de RL
são dimensionados proporcionalmente a $Q_\text{ref}$ e $s_\text{ref}$.

\subsection{Configuração de Custos e Protocolo de Comparação}

Os parâmetros de custo são uniformes para todas as lojas e políticas:
$h = 1.0$, $c_s = 5.0$, $c_o^{fix} = 50.0$ e $c_o^{var} = 0.5$.
Todas as políticas são avaliadas sobre a mesma série de demanda
real, no mesmo horizonte fixo de $T = 38$ ciclos, com calibração
individualizada por loja conforme descrito acima. Os cinco
indicadores definidos na Seção~3.3 são computados ao fim de cada
simulação para cada combinação (política, loja).

% =========================================================
\section{Resultados e Discussão}
% =========================================================

Esta seção apresenta a comparação global das doze políticas, a
análise estratificada por grupo de demanda, o comportamento do
efeito \textit{bullwhip}, uma discussão dos mecanismos subjacentes
e a comparação direta com \citet{zabraoui:2025}.

\subsection{Comparação Global entre Políticas}

A Tabela~\ref{tab:global} apresenta as médias dos cinco indicadores
sobre as quinze lojas, ordenadas por CTI crescente. Os melhores
valores em cada coluna estão destacados em negrito.

\begin{table}[ht]
\centering
\small
\caption{Médias dos indicadores de desempenho por política (15 lojas).
Melhores valores por coluna em \textbf{negrito}. $^\dagger$SARSA e PPO
apresentam degeneração ($FP = 0$ e $FP = 0.98$, respectivamente);
resultados degenerados não constituem soluções operacionalmente viáveis.}
\label{tab:global}
\begin{tabular}{@{}lrrrrr@{}}
\toprule
Política &
\makecell[r]{CTI\\(R\$)} &
\makecell[r]{NS} &
\makecell[r]{TR} &
\makecell[r]{BE} &
\makecell[r]{FP} \\
\midrule
SARSA$^\dagger$              & \textbf{16.139} & 0{,}785 & 0{,}200 & \textbf{0{,}00} & \textbf{0{,}00} \\
Jornaleiro (\textit{Newsvendor}) & 16.402 & 0{,}794 & 0{,}111 & 0{,}003 & 0{,}27 \\
Rede Q Profunda (DQN)        & 17.974 & 0{,}806 & 0{,}078 & 4{,}1   & 0{,}37 \\
GA-DQN (Híbrido)             & 18.502 & 0{,}821 & 0{,}050 & 5{,}0   & 0{,}31 \\
Enxame de Partículas (PSO)   & 19.100 & 0{,}819 & 0{,}056 & 6{,}2   & 0{,}31 \\
GA-PPO (Híbrido)             & 19.801 & 0{,}833 & 0{,}050 & 6{,}6   & 0{,}35 \\
Algoritmo Genético (GA)      & 20.015 & 0{,}825 & 0{,}061 & 19{,}7  & 0{,}30 \\
Evolução Diferencial (ED)    & 20.252 & 0{,}813 & 0{,}056 & 31{,}0  & 0{,}27 \\
\textit{Sim. Annealing} (SA) & 20.967 & 0{,}820 & 0{,}056 & 36{,}4  & 0{,}27 \\
Lote Econômico (EOQ)         & 31.034 & 0{,}840 & 0{,}050 & 23{,}0  & 0{,}20 \\
Política PPO$^\dagger$       & 34.637 & \textbf{0{,}848} & \textbf{0{,}039} & 15{,}9 & 0{,}98 \\
Política $(s,S)$             & 35.589 & 0{,}843 & 0{,}050 & 119{,}4 & 0{,}15 \\
\bottomrule
\end{tabular}
\end{table}

Os dois extremos revelam os riscos da otimização por única métrica:
o SARSA$^\dagger$ converge à estratégia de não emitir pedidos
(FP\,=\,0), obtendo o menor CTI ao preço da maior taxa de ruptura
($0.200$); o PPO$^\dagger$ emite pedidos em $98\%$ dos ciclos,
atingindo o maior NS mas incorrendo em CTI inviável de
R\$\,34.637. Entre as políticas práticas (sem degeneração), o
GA-DQN apresentou o equilíbrio mais favorável entre as políticas
práticas, com CTI apenas $13\%$ acima do Jornaleiro e TR $55\%$
menor. A política~$(s,S)$, frequentemente utilizada no varejo
regional, tem CTI $92\%$ superior ao GA-DQN e efeito
\textit{bullwhip} de 119 vezes, indicando limitações consistentes
nesse regime de demanda.

Ressalva-se que os resultados acima comparam médias sobre as quinze
lojas sem testes de significância estatística formal. A ausência de
réplicas independentes impede a aplicação de testes como o Wilcoxon
\textit{signed-rank} ou Friedman com \textit{post-hoc} de Nemenyi;
diferenças pequenas entre políticas próximas (como GA-DQN, GA-PPO
e PSO) devem ser interpretadas com cautela. Essa limitação é
discutida na Seção~8.3.

\subsection{Análise Estratificada por Grupo de Demanda}

A estratificação por grupo evidencia que não existe política
universalmente dominante, resultado de maior relevância prática
para gestores de \textit{supply chain}.

Grupo I ($\hat{\mu} > 500$, cinco lojas): estoque inicial
adequado garante NS\,=\,$1.00$ para quase todas as políticas; o
diferenciador é o custo. O EOQ gasta em média R\$\,85.956 por loja
contra R\$\,45.739 das políticas de menor frequência de pedidos,
redução de $47\%$ apenas pela troca de heurística.

Grupo II ($50 < \hat{\mu} \leq 500$, seis lojas): estrato mais
discriminante e de maior relevância prática. O GA-DQN destaca-se
com NS\,=\,$0.933$ e CTI\,=\,R\$\,10.854, a melhor combinação
custo-serviço do estudo, atribuível ao \textit{replay buffer} de
alta qualidade gerado pelo GA~\citep{oroojlooyjadid:2022}.

Grupo III ($\hat{\mu} \leq 50$, quatro lojas): três lojas
não superam NS de $0.42$ com nenhuma política. O CV próximo de 3
cria uma barreira estrutural em que pedidos compatíveis com a média
não cobrem picos e pedidos elevados aumentam o custo sem ganho de
serviço. Nesse estrato, o Jornaleiro é a política mais econômica.

\subsection{Estabilidade Logística e Efeito Bullwhip}

O efeito \textit{bullwhip}~\citep{lee:1997} mede quanto a
variabilidade dos pedidos supera a da demanda; valores elevados
implicam oscilações desnecessárias e custos logísticos crescentes
a montante. Os resultados revelam amplitude expressiva: excluindo
o SARSA (degenerado, BE\,=\,0), as políticas práticas variam de
$0.003$ (Jornaleiro) a $119.4$ em~$(s,S)$, razão superior a 40.000
vezes. Entre as políticas práticas, DQN ($4.1$) e GA-DQN ($5.0$)
apresentam menor amplificação por aprenderem a antecipar picos de
demanda; PSO ($6.2$) é a melhor meta-heurística nessa dimensão,
enquanto SA ($36.4$) e ED ($31.0$) ficam elevados por otimizarem
$(s,Q)$ pela média histórica. Este indicador, avaliado aqui como
critério primário para todas as doze políticas, permite analisar o
impacto de cada abordagem sobre a estabilidade de toda a cadeia,
dimensão crítica para distribuidores e operadores logísticos.

A Figura~\ref{fig:barras} sintetiza as três métricas primárias.

\begin{figure}[ht]
  \centering
  \includegraphics[width=\linewidth]{figuras/pb_barras.png}
  \caption{Desempenho comparativo das doze políticas segundo CTI
  médio, NS médio e BE médio (barras ordenadas do melhor para o pior
  em cada indicador).}
  \label{fig:barras}
\end{figure}

\subsection{Discussão: Mecanismos e Comportamentos Observados}

Os resultados permitem identificar quatro padrões explicativos.

\textbf{Por que o GA-DQN supera o DQN isolado.} O GA pré-popula o
\textit{replay buffer} com 5.000 avaliações simuladas ($100 \times 50$
gerações), reduzindo a dependência de exploração em horizonte de
apenas 570 transições reais ($15 \times 38$) e contornando o
requisito de da ordem de $10^6$ passos do DQN~\citep{mnih:2015}.

\textbf{Por que o PPO converge para política de alta frequência de pedidos.}
Sem inicialização, o gradiente da política estocástica converge para
maximizar NS imediato, ignorando o custo de manutenção
(FP\,=\,$0.98$, CTI\,=\,R\$\,34.637). O \textit{buffer}
pré-populado pelo GA força o agente a partir de uma região de
equilíbrio custo-serviço, evitando essa degeneração no GA-PPO.

\textbf{Por que o SARSA apresenta convergência inadequada.}
O SARSA tabular com espaço de estados $5\times10$ tem
expressividade insuficiente para o regime intermitente: com demanda
nula na maior parte dos ciclos, a política \textit{on-policy}
tende a aprender que não emitir pedido é a ação de menor custo
imediato, convergindo para FP\,=\,0 e TR\,=\,0.200.

\textbf{Meta-heurísticas versus RL em horizonte curto.}
Meta-heurísticas otimizam diretamente $(s,Q)$ pela simulação do
horizonte completo, integrando toda a informação disponível de uma
vez; agentes de RL aprendem passo a passo, exigindo muito mais
episódios para o mesmo padrão. Em $T = 38$, a otimização global é
estruturalmente mais adequada, justificando o papel do GA como
inicializador nas arquiteturas híbridas.

\textbf{Custo computacional relativo.} Políticas clássicas são
analíticas (tempo constante); meta-heurísticas requerem segundos a
minutos por loja, mas produzem parâmetros calculados uma única vez;
agentes de RL exigem centenas de episódios de treinamento.
Arquiteturas híbridas somam ambos os custos, porém o treinamento é
realizado uma única vez por loja, tornando-as viáveis como rotina
periódica em sistemas ERP.

\subsection{Comparação com a Literatura}

Considerando os cinco indicadores conjuntamente, o GA-DQN
apresentou o equilíbrio mais favorável entre as políticas avaliadas:
CTI de R\$\,18.502 (apenas
$13\%$ acima do Jornaleiro, que tem TR mais que o dobro), NS de
$0.821$ equiparável ao das meta-heurísticas isoladas
($0.813$--$0.825$) com CTI menor, BE de $5.0$
e TR de $0.050$ equivalente à~$(s,S)$ porém com custo $92\%$ menor.

Esses resultados confirmam, em contexto de $CV$ entre $2.5$ e $5.4$,
a vantagem do GA-DQN reportada por \citet{zabraoui:2025} para o
Walmart~M5 (CV típico $< 1$), sugerindo que a arquitetura apresenta desempenho consistente
entre regimes de demanda e ampliando seu escopo de aplicação.

% =========================================================
\section{Conclusão}
% =========================================================

\subsection{Principais Resultados}

A comparação sistemática de doze políticas em dados operacionais
proprietários de varejo regional brasileiro com $CV$ entre $2.5$
e $5.4$ sugere que a arquitetura híbrida GA-DQN apresentou o
equilíbrio mais favorável entre custo total, nível de serviço e
estabilidade logística entre as políticas avaliadas. A hipótese
central, de que arquiteturas híbridas GA-RL apresentam equilíbrio
mais favorável do que políticas clássicas e agentes de RL isolados,
é apoiada pelos resultados observados. A política~$(s,S)$,
frequentemente utilizada no varejo regional brasileiro, apresentou
desempenho consistentemente inferior: superada em custo e
estabilidade por todas as alternativas práticas avaliadas, com
efeito \textit{bullwhip} aproximadamente 24 vezes maior que o
GA-DQN e CTI $92\%$ superior.

\subsection{Implicações Práticas}

Com base na análise estratificada, três recomendações emergem
diretamente dos resultados:

\begin{itemize}
  \item Adotar GA-DQN como política padrão para o Grupo II
    (lojas de volume médio): NS\,=\,$0.933$ e
    CTI\,=\,R\$\,10.854 justificam o custo computacional adicional
    frente às heurísticas clássicas. O treinamento é realizado uma
    única vez por loja e produto.
  \item Considerar a substituição da política~$(s,S)$:
    consistentemente superada em custo e estabilidade logística
    pelas demais alternativas práticas em todos os estratos
    avaliados nesse regime de demanda. A migração
    direta para o Jornaleiro, sem custo computacional adicional,
    já contribui para reduzir o problema de amplificação da cadeia.
  \item Utilizar o Jornaleiro no Grupo~III (baixo volume): sua
    simplicidade e baixo CTI superam abordagens sofisticadas nesse
    estrato, onde a barreira estrutural de NS limita qualquer
    política independentemente de sua complexidade.
\end{itemize}

Para \emph{redes de varejo regional}, a adoção do GA-DQN pode
representar redução da ordem de $48\%$ no CTI frente à~$(s,S)$
(diferença de R\$\,17.087 por loja por período, sujeita às premissas
do estudo), sem perda de nível de serviço.
Para \emph{distribuidores e operadores logísticos}, a redução do
efeito \textit{bullwhip} de $119$ para $5$ estabiliza a cadeia a
montante, diminuindo custos de transporte e armazenagem.
Para \emph{pequenas e médias empresas} sem equipe de
\textit{data science}, a arquitetura requer apenas histórico de
vendas disponível em qualquer sistema ERP e pode ser
operacionalizada como rotina periódica automatizada.
Do ponto de vista científico, o estudo oferece uma base de
comparação reproduzível em regime de alta intermitência,
replicável para qualquer produto com CV~$>1.5$.

\subsection{Limitações}

O estudo apresenta quatro limitações que devem ser consideradas na
interpretação dos resultados. Primeiro, adota simulação
determinística com série de demanda estática e sem réplicas
independentes, o que pode subestimar a variância dos estimadores
de desempenho. Segundo, restringe-se a um único SKU (produto 48130)
em uma única região (Paraíba), limitando a generalização direta
para outros produtos ou mercados. Terceiro, não considera restrições
de capacidade de armazenagem nem prazo de entrega estocástico,
presentes em ambientes reais de maior complexidade. Quarto, a
comparação entre políticas baseia-se em médias sem testes de
significância estatística formal; a ausência de réplicas independentes
impede a aplicação de procedimentos como o Wilcoxon
\textit{signed-rank} ou o Friedman com \textit{post-hoc} de Nemenyi,
de modo que diferenças pequenas entre políticas próximas
(particularmente GA-DQN, GA-PPO e PSO) podem não ser
estatisticamente significativas. Quinto, os resultados representam
desempenho \textit{in-sample}: a mesma série histórica é utilizada
para calibração e avaliação das políticas, o que pode favorecer
abordagens com maior capacidade de ajuste aos dados observados;
validação \textit{out-of-sample} com janela deslizante é necessária
para confirmar a generalização.

\subsection{Trabalhos Futuros}

Como direções futuras, propõe-se: avaliação \textit{out-of-sample}
com janela deslizante para validar a generalização; extensão a
múltiplos produtos com restrições de capacidade de armazenagem;
incorporação de prazo de entrega estocástico; integração com modelos
de previsão de demanda (Croston, \textit{Long Short-Term Memory}
(LSTM)) para alimentar as decisões com estimativas prospectivas;
e exploração de arquiteturas de RL multiagente para coordenação
explícita entre lojas.

\bibliographystyle{sbpo}
\bibliography{exemplo-latex}

\end{document}
